% !TeX spellcheck = ru_RU
% !TEX root = vkr.tex

% расписать нормальное завершение с указанымы решенными задачами

\section{Заключение}

  В ходе выполнения выпускной квалификационной работы получены следующие результаты.
  \begin{itemize}
    \item Выполнен обзор открытых процессорных архитектур (RISC‑V, MIPS, ARM), легковесных UNIX‑подобных ОС (xv6, Minix, uClinux) и компактных С-компиляторов C (TinyCC, LCC). В качестве технологического стека выбраны RISC‑V/ xv6/ TinyCC.
    \item Спроектирована модульная архитектура комплекса, описывающая роли и интерфейсы виртуального процессора, подсистем памяти и периферии, ОС, компилятора и средств трассировки.
    \item Разработано синтезируемое бесконвейерное ядро RISC‑V RV32I с помощью SystemVerilog с поддержкой базового набора инструкций и генерацией VCD‑файлов. Создана моделируемая вычислительная платформа, включающая память, простейшую периферию (UART, таймер) и загрузчик, в двух конфигурациях: run\_c\_version для выполнения «голого» С‑кода и xv6\_version для запуска ОС xv6.
    \item Выполнена адаптация и сборка UNIX‑подобной ОС xv6 для созданного процессорного ядра (модифицированная версия, не требующая MMU и S‑режима), обеспечена корректная загрузка и выполнение ядра на виртуальном процессоре.
    \item Адаптирован легковесный компилятор TinyCC для работы в среде xv6/RISC-V; выполнена его интеграция в образ ОС, что позволяет компилировать и запускать C-программы на платформе без участия хостовой операционной системы.
    \item Работоспособность комплекса подтверждена успешным выполнением минимальной ОС xv6 и пользовательских программ на виртуальном процессоре; верифицирована корректность управления процессами, памятью и вводом‑выводом. Воспроизводимость платформы обеспечена унифицированными репозиториями (RTL‑модель, C++ ISA‑симулятор, Docker‑образ, тестовые сценарии).
    \item На основе предложенного программного комплекса разработаны лекционные материалы, сценарии симуляции и подготовлена подробная документация, описывающая архитектуру и структуру программного комплекса (компонентный состав, инструкции по сборке и методика использования в учебном процессе).
  \end{itemize}
	код проекта представлен в репозитории \href{https://github.com/KiselkovD/vcpu-for-xv6-rv32i-edu-platform}{KiselkovD/vcpu-for-xv6-rv32i-edu-platform}